\section{State, Frame, and Descriptive Structure}

\begin{definition}[Behavioural state]
A \textbf{behavioural state} $s \in \mathcal{S}$ is a sustained mode of action
together with its associated posture, attention pattern, local environment, and
reward expectation.
\end{definition}

\begin{remark}
This is a modelling definition, not an empirical claim that every behavioural
state is a sharply separable biological object. Its purpose is to let the book
talk about persistence and transition without collapsing everything into the
single word ``motivation.''
\end{remark}

\begin{definition}[State profile]
For descriptive purposes, a behavioural state may be assigned a
\textbf{state profile}
\[
  \Pi(s) = \bigl(\Reward(s), \Switch(s), \Env(s)\bigr),
\]
where:
\begin{itemize}[nosep]
  \item $\Reward(s)$ denotes the immediate reward pull of remaining in $s$,
  \item $\Switch(s)$ denotes the cognitive cost of loading a competing task
        frame,
  \item $\Env(s)$ denotes environmental reinforcement of the current state.
\end{itemize}
\end{definition}

\begin{discussion}
The tuple $\Pi(s)$ is more defensible than prematurely forcing a universal
scalar. It says only that these three components are analytically useful and
should be named separately. That separation is already valuable for later case
analysis and intervention design, because it identifies different control
levers: reward structure, context loading, and environmental coupling.
\end{discussion}

\begin{remark}
At this stage, $\Pi(s)$ should be read as an \textbf{ordered descriptive
profile}, not as a measured coordinate system with fixed units. The chapter is
claiming that these dimensions are useful to separate, not that they have
already been standardized across persons, tasks, or laboratories.
\end{remark}
